% Generated by escalation/run_bench_script/make_figures_tex.py -- do not edit.
% The caption text lives in that chart's plot script, next to the numbers it
% describes; re-run the script to update this file.
\begin{figure}[p]
\centering
\begin{subfigure}{\linewidth}
  \centering
  \panelplot{\plotdir/manager_vs_single_call_across_model_scale_lcb-100_1_pass_128k_max_tokens_reasoning_on_light.png}
  \caption{Accuracy by model scale.}
\end{subfigure}\\[0.7ex]
\begin{subfigure}{\linewidth}
  \centering
  \panelplot{\plotdir/manager_vs_single_call_across_model_scale_lcb-100_1_pass_128k_max_tokens_reasoning_on_bars_light.png}
  \caption{The same numbers as paired bars.}
\end{subfigure}
\caption{\textbf{Manager vs. single, reasoning on, 128k, one pass.} \footnotesize\normalfont
Fill: light = single call (one call, no tools), dark = with manager; in (a) a half-and-half marker means both arms scored the same. 128k max tokens, reasoning ON - effort:high for Kimi and Opus, and a 20k reasoning budget requested for Qwen and Minimax that providers often did not honour. $\Delta$ = with manager $-$ single call, in percentage points. One pass means no within-model repeat, so no per-model p is shown. Non-empty completions per 100 (single$\to$manager): Qwen3.6-35B 34$\to$68, Minimax-M3 92$\to$93, Kimi-K3 97$\to$100, Opus-5 97$\to$100. Empty or truncated output scores as a fail, so $\Delta$ partly tracks emit rate. The curves in (a) connect all four models (monotone cubic): they interpolate, they are not fits. Each block heads with the model's size; Opus-5's is not public, so it heads with its name, sits at a placeholder x in (a) and last in (b). Tukey groups: models sharing a letter have indistinguishable $\Delta$ (HSD, $\alpha$ 0.05); HSD ignores that all models saw the same 100 problems, so it is conservative. Pairwise Tukey p: 35B vs 428B 0.19 n.s.; 35B vs 2.8T 0.0088; 35B vs Opus-5 0.19 n.s.; 428B vs 2.8T 0.65 n.s.; 428B vs Opus-5 1 n.s.; 2.8T vs Opus-5 0.65 n.s.}
\end{figure}
